\begin{abstract}
Large-scale unsupervised language models (LMs) acquire general world knowledge and certain reasoning capabilities, yet fine-grained regulation of their outputs remains challenging due to the fully unsupervised nature of their training process. To address this, existing approaches introduce human-generated preference data that rank model outputs, subsequently refining the original unsupervised LM so that it better conforms to these preferences—typically using reinforcement learning from human feedback (RLHF). RLHF, however, is both intricate and frequently unstable, involving a two-stage procedure: first, training a reward model to approximate human judgments, and second, fine-tuning the LM with reinforcement learning to maximize the inferred reward while avoiding excessive deviation from the base model. This work proposes a novel reward model parameterization for RLHF that makes it possible to recover the optimal policy in closed form, reducing the RLHF framework to a straightforward classification loss. We introduce an algorithm based on this insight, termed \textit{Direct Preference Optimization} (DPO), which is stable, efficient, and resource-friendly; it does not require model sampling during the fine-tuning phase nor extensive hyperparameter searches. Empirically, DPO is able to adapt LMs to human preferences on par with or superior to current techniques. Remarkably, DPO outperforms RLHF methods based on PPO when guiding generation sentiment, and it achieves equal or superior performance in tasks like summarization and single-turn dialogue, while offering considerably reduced complexity in implementation and training.
\end{abstract}

\section{Introduction}
Large unsupervised language models (LMs) trained on massive datasets have demonstrated remarkable emergent abilities~\citep{chowdhery2022palm, brown2020language, touvron2023llama,bubeck2023sparks}. Nevertheless, the underlying training data reflects a broad spectrum of human intentions, priorities, and expertise. Many of these human-driven goals and competencies are not always ideal for replication; for instance, though we want an AI assistant for programming to \textit{recognize} frequent coding mistakes for the purpose of correction, we would prefer its code generation to reflect the infrequent but superior-quality coding instances present in the data. Likewise, while a language model should \textit{comprehend} prevalent misconceptions (such as those accepted by half the population), it should not endorse such misunderstandings in half its corresponding outputs. Accordingly, discerning and promoting the \emph{appropriate behaviors and responses} from among the model’s extensive \textit{knowledge and capabilities} is essential to the development of AI systems that are reliable, effective, and amenable to control \citep{ouyang2022training}. Current methodologies commonly use reinforcement learning (RL) to align LMs with human values, but in this work, we demonstrate that the objective underpinning RL-based preference alignment can be satisfied directly by a binary cross-entropy loss, thereby streamlining the preference optimization workflow.

Broadly speaking, prevalent approaches imbue language models with target behaviors by leveraging curated datasets of human preferences, which encode safe and useful conduct. This preference learning process is introduced after the unsupervised pre-training phase on extensive text corpora. Although supervised fine-tuning on exemplary human-generated responses is a direct strategy, reinforcement learning from human (or AI-generated) feedback (RLHF/RLAIF; \citep{christiano2017deep,bai2022constitutional}) has emerged as the most effective paradigm. RLHF frameworks involve training a reward model on human preference data, which is subsequently used to guide the LM policy via RL towards producing outputs that receive high rewards, all while mitigating significant deviation from the pre-trained baseline. While RLHF leads to models that excel at conversation and programming, the associated process is substantially more intricate than standard supervised learning—entailing the training of multiple models and continuous policy sampling within the training loop, which substantially increases computational requirements.

This work presents a method for optimizing language models in accordance with human preferences without requiring explicit reward modeling or the use of reinforcement learning. We introduce \textit{{Direct Preference Optimization} (DPO)}, a technique that effectively pursues the same objective as traditional RLHF algorithms (maximizing reward under a KL-divergence penalty), but offers a more accessible and easier-to-train solution. The DPO procedure raises the log odds of preferred responses relative to less-preferred ones, introducing a data-dependent, sample-wise weighting that mitigates the model degradation often seen when employing a simple probability ratio approach. Similar to prior approaches, DPO employs a formal preference model (e.g., the Bradley-Terry model; \cite{bradley1952rankanalysis}) to gauge the alignment between reward functions and observed preference judgments. However, while previous frameworks typically leverage this model to craft a loss for training a reward model, and then optimize a policy against the learned reward, DPO applies a reparameterization to define the loss directly in terms of the policy. Provided a collection of preference annotations over model outputs, DPO enables policy optimization through a straightforward binary cross entropy loss, resulting in the optimal policy relative to an implicitly inferred reward function that matches the observed preferences.

The central advance of our study is the proposal of Direct Preference Optimization (DPO), a robust algorithm for training language models based on preferences, which entirely avoids RL. Experimental results demonstrate that DPO achieves performance on par with, or superior to, existing methods—including those based on PPO—for tasks such as sentiment control, summarization, and conversational modeling, utilizing language models containing as many as 6B parameters.

\section{Related Work}

Progressively larger self-supervised language models have demonstrated the ability to address certain tasks without task-specific training, employing zero-shot approaches \citep{radford2019language} or relying on few-shot prompting strategies \citep{gpt3,megatron,chowdhery2022palm}. Nonetheless, substantial gains in performance on downstream applications and better adherence to user intent are realized when these models undergo additional fine-tuning with curated datasets comprising instructions paired with human-authored completions \citep{mishra-etal-2022-cross,sanh2022multitask,chung2022scaling,thoppilan2022lamda}. This process, referred to as `instruction-tuning', enhances LLMs’ capacity to handle novel instructions outside the scope of their training data, thereby broadening their practical utility \citep{chung2022scaling}. Although instruction tuning yields significant improvements, collecting \textit{relative} human assessments of response quality tends to be more scalable compared to sourcing expert-annotated demonstrations. Consequently, subsequent research has prioritized fine-tuning LLMs using collections of human preference data, which has resulted in better outcomes for translation \citep{kreutzer-etal-2018-reliability}, summarization \citep{stiennon2022learning,ziegler2020finetuning}, narrative generation \citep{ziegler2020finetuning}, and instruction adherence \citep{ouyang2022training,ramamurthy2023is}. The prevailing approach entails first training a neural reward model that reflects the dataset of human preferences, typically within a framework such as the Bradley-Terry model \citep{bradley1952rankanalysis}; the language model is then adjusted to optimize this reward via reinforcement learning techniques, often employing REINFORCE \citep{williams1992reinforce}, proximal policy optimization (PPO; \cite{schulman2017proximal}), or related algorithms \citep{ramamurthy2023is}. In a related vein, some studies utilize LLMs fine-tuned through human feedback to synthesize additional preference datasets that focus on targeted features like safety or non-harmfulness \citep{bai2022constitutional}, relying on weak human supervision delivered as text-based guidelines for LLM annotation. Collectively, these methodologies reflect the integration of two research streams: the use of reinforcement learning to train language models for diverse tasks~\citep{Ranzato2015SequenceLT,paulus2018a,wu2018learning}, and general frameworks for inferring from human preferences \citep{christiano2017deep,kupcsik2018learning}. While optimizing using relative preferences is attractive, applying reinforcement learning to fine-tune large-scale language models is still a significant technical hurdle; this paper introduces a theoretically sound method for preference optimization that circumvents the need for RL.

Beyond language-specific domains, the study of learning policies from preference data has been extensively explored within both bandit and reinforcement learning frameworks, yielding a variety of methodologies. In contextual bandit settings, where feedback is in the form of preferences or rankings over actions rather than numerical rewards, the problem is formalized as the contextual dueling bandit (CDB; \cite{yue2012karmed,dudik2015contextual}). Without access to absolute reward signals, theoretical work on CDBs replaces the optimal policy concept with that of a \textit{von Neumann winner}—a policy guaranteed to achieve an expected win rate of at least 50\% against any competing policy \citep{dudik2015contextual}. A key distinction of CDBs is the provision of preference feedback in an online manner; by contrast, human preference-based learning generally relies on a static, finite dataset of previously annotated action pairs \citep{yan2022human}. 

Likewise, in \textit{preference-based RL} (PbRL), policy improvement is guided by binary preferences derived from an implicit, unknown 'scoring' function, instead of direct reward feedback \citep{BusaFekete2014,ruiz2023dueling}. Approaches in this area range from algorithms that can leverage off-policy preference observations to those that involve reconstructing the latent scoring (reward) function and subsequently performing policy optimization with respect to it \citep{jain2013learning,BusaFekete2014,christiano2017deep,sadigh2017active,kupcsik2018learning}. In contrast to these two-step strategies, we introduce a unified framework that directly learns a policy optimized for preference satisfaction in a single stage.

\section{Preliminaries}\label{section:prelims}

We begin by summarizing the RLHF workflow as originally articulated by \citeauthor{ziegler2020finetuning} and subsequently refined in \citep{stiennon2022learning, bai2022training, ouyang2022training}. The pipeline typically comprises three core steps: (1) supervised fine-tuning (SFT), (2) collection of preferences and construction of a reward model, and (3) reinforcement learning optimization.

\textbf{SFT}: The process starts with applying supervised fine-tuning to a pre-trained language model, using curated data that aligns with the intended end task (such as dialogue or summarization). This yields a task-specialized model $\pi^\text{SFT}$.

\textbf{Reward Modelling Phase}: The subsequent phase involves leveraging $\pi^\text{SFT}$ to generate, given input prompts $x$, pairs of candidate responses $(y_1, y_2)\sim \pi^\text{SFT}(y \mid x)$. Human annotators are then tasked with expressing a preference between these two, denoted as $y_w\succ y_l \mid x$, where $y_w$ and $y_l$ represent, respectively, the chosen and rejected outputs. Preferences are postulated to be governed by a hidden reward function $r^*(y, x)$, which remains unobserved. Several modeling techniques are employed to characterize these preferences, with the Bradley-Terry (BT) \cite{bradley1952rankanalysis} model frequently adopted. More sophisticated models, such as Plackett-Luce ranking models \citep{plackett1975analysis, luce2012individual}, may also be applicable when dealing with multiple ranked completions. According to the BT framework, the probability distribution over human preferences $p^*$ is given by:

\begin{equation}\label{eq:bradley-terry}
    p^*(y_1\succ y_2 \mid x)=\frac{\exp\left(r^*(x, y_1)\right)}{\exp\left(r^*(x, y_1)\right) + \exp\left(r^*(x, y_2)\right)}.
\end{equation}

Given a fixed dataset of pairwise preferences $\mathcal{D}=\bigl\{x^{(i)}, y_w^{(i)}, y_l^{(i)}\bigr\}_{i=1}^N$ drawn according to $p^*$, one may introduce a parameterized reward function $r_{\phi}(x, y)$ and fit its parameters through maximum likelihood estimation. By rephrasing this as a binary classification task, the following negative log-likelihood loss is obtained:

\begin{equation}\label{eq:reward_model}
    \mathcal{L}_R(r_{\phi}, \mathcal{D}) = -\mathbb{E}_{(x, y_w, y_l)\sim \mathcal{D}}\bigl[\log \sigma(r_{\phi}(x, y_w)- r_{\phi}(x, y_l))\bigr]
\end{equation}

where $\sigma$ denotes the logistic sigmoid function. For language models, $r_{\phi}(x, y)$ is generally built by initializing from an SFT model $\pi^\text{SFT}(y \mid x)$ and attaching an additional linear layer on top of the final transformer representation to produce a scalar reward score~\cite{ziegler2020finetuning}. To stabilize the learned reward, normalization is often employed, enforcing  $\mathbb{E}_{x,y\sim \mathcal{D}}\left[r_\phi(x, y)\right] = 0$ for each $x$, as in previous studies.

\textbf{RL Fine-Tuning Phase}: In the reinforcement learning stage, the inferred reward function serves as the supervisory signal for the language model's outputs. Building on prior literature~\citep{jaques2017sequence, jaques2020human}, the optimization target is framed as:

\begin{equation}\label{eq:RL}
\max_{\pi_{\theta}}  \mathbb{E}_{x\sim \mathcal{D}, y\sim \pi_{\theta}(y \mid x)}\bigl[r_{\phi}(x, y)\bigr] - \beta\mathbb{D}_{\textrm{KL}}\bigl[\pi_{\theta}(y\mid x)\mid \mid \pi_\text{ref}(y\mid x)\bigr],
\end{equation}

Here, $\beta$ regulates the amount by which $\pi_\theta$ is allowed to diverge from the reference distribution $\pi_\text{ref}$—namely, the starting SFT model $\pi^\text{SFT}$. The policy $\pi_\theta$ is typically initialized from $\pi^\text{SFT}$ as well. The KL penalty plays a key role by constraining the model close to domains where reward estimates remain reliable and also discourages collapse toward only a few highly-rewarded responses, thus preserving response diversity. Given the non-differentiability of the sampling process in language, such an objective necessitates reinforcement learning algorithms. Prevailing methodologies~\citep{ziegler2020finetuning, stiennon2022learning, bai2022training, ouyang2022training} reformulate the reward as ${r(x, y) = r_{\phi}(x, y) -\beta (\log \pi_{\theta}(y\mid x) - \log \pi_\text{ref}(y\mid x))}$ and employ PPO~\cite{schulman2017proximal} to perform the optimization.

\section{Direct Preference Optimization}\label{sec:DPO}

Recognizing the limitations of conventional reinforcement learning when scaling up to fine-tune large language models, we seek to formulate a more direct preference-based policy optimization scheme. In contrast to standard RLHF pipelines that first learn an explicit reward estimator and subsequently train the policy with RL, our method introduces a strategic parameterization of the reward such that the corresponding optimal policy can be recovered analytically—eliminating the need for reinforcement learning altogether. As detailed below, this approach builds upon a principled mapping between rewards and optimal policies, allowing us to substitute a policy-based loss for the reward-model loss. This change-of-variable method sidesteps separate reward-model fitting, while remaining compatible with prevalent human preference models such as the Bradley-Terry framework. In this formulation, the policy network effectively embodies both the language generation mechanism and the implied

\textbf{DPO objective derivation.} Our starting point aligns with earlier works and employs the RL objective in Eq.~\ref{eq:RL} with a general reward function $r$. Building on results from previous literature~\citep{peters2007reinforcement, peng2019advantage, korbak2022reinforcement, go2023aligning}, it is well known that the solution optimizing the KL-regularized reward maximization in Eq.~\ref{eq:RL} adopts the following structure:

\begin{equation}\label{eq:op_policy}
    \pi_r(y\mid x) = \frac{1}{Z(x)}\pi_\text{ref}(y\mid x)\exp\left(\frac{1}{\beta}r(x, y)\right),
\end{equation}

where $Z(x) =\sum_{y}\pi_\text{ref}(y\mid x)\exp\left(\frac{1}{\beta}r(x, y)\right)$ denotes the normalizing constant (partition function). For details of the derivation, see Appendix \ref{app:derivation1}. However, even substituting in the MLE estimate $r_{\phi}$ for the underlying reward $r^*$ does not resolve the computational difficulty associated with estimating $Z(x)$~\citep{korbak2022reinforcement, go2023aligning}, thereby making practical application nontrivial. Nevertheless, by manipulating Eq.~\ref{eq:op_policy}, the reward $r$ can be reformulated in terms of the optimal policy $\pi_r$, the reference policy $\pi_\text{ref}$, and $Z(\cdot)$. Taking logarithms on both sides and rearranging terms yields:

\begin{equation}\label{eq:main_eq}
    r(x,y) =\beta \log \frac{\pi_r(y\mid x)}{\pi_\text{ref}(y\mid x)} + \beta \log Z(x).
\end{equation}

Applying this transformation to both the ground-truth reward $r^*$ and its associated optimal policy $\pi^*$ is direct. Critically, the Bradley-Terry model utilizes only the difference in rewards for two outputs, i.e., ${p^*(y_1 \succ y_2 \mid x) = \sigma(r^*(x, y_1) - r^*(x, y_2))}$. By plugging the representation from Eq.~\ref{eq:main_eq} for $r^*(x,y)$ into the preference probability given by Eq.~\ref{eq:bradley-terry}, the normalization term cancels out, and the probability can be expressed solely using $\pi^*$ and $\pi_\text{ref}$. Accordingly, the policy $\pi^*$ learned by RLHF under the Bradley-Terry framework meets the relation:

\begin{equation}\label{eq:objective}
    p^*(y_1\succ y_2 \mid x)=\frac{1}{1 + \exp\left(\beta \log \frac{\pi^*(y_2\mid x)}{\pi_\text{ref}(y_2\mid x)} - \beta \log \frac{\pi^*(y_1\mid x)}{\pi_\text{ref}(y_1\mid x)}\right)}
\end{equation}

Appendix~\ref{app:derivation2} details the derivation. Although Eq.~\ref{eq:objective} is based on the Bradley-Terry construction, the same reasoning extends to more general preference aggregation models, such as Plackett-Luce~\citep{plackett1975analysis, luce2012individual}, as discussed in Appendix~\ref{app:plackett_luce_models}.

Given this reformulation, we now relate human preference probabilities directly to the optimal policy, circumventing the explicit reward model. This makes it natural to define a maximum likelihood estimation problem for a learnable policy $\pi_\theta$. Mirroring the strategy used in reward modeling (cf. Eq.~\ref{eq:reward_model}), the policy learning objective becomes:

\begin{equation}\label{eq:optimum_model}
    \mathcal{L}_\text{DPO}(\pi_{\theta}; \pi_\text{ref}) = -\mathbb{E}_{(x, y_w, y_l)\sim \mathcal{D}}\left[\log \sigma \left(\beta \log \frac{\pi_{\theta}(y_w\mid x)}{\pi_\text{ref}(y_w\mid x)} - \beta \log \frac{\pi_{\theta}(y_l\mid x)}{\pi_\

In this formulation, we represent the reward implicitly via an alternative parametrization, whereby the resulting optimal policy aligns with $\pi_\theta$. Notably, our approach corresponds to fitting a reparameterized Bradley-Terry model, granting us desirable theoretical guarantees such as consistency under appropriate assumptions regarding the distribution of preference data \cite{bong2022generalized}. Section~\ref{sec:theory} provides additional discussion on the theoretical characteristics of DPO and how they relate to previous literature.

\textbf{Mechanistic interpretation of the DPO update:} To gain insight into how DPO operates, one can examine the gradient of its loss function $\mathcal{L}_\text{DPO}$. Specifically, the gradient with respect to $\theta$ is given by:

\begin{multline*}\label{eq:gradient}
    \nabla_\theta \mathcal{L}_\text{DPO}(\pi_\theta;\pi_\text{ref}) = \\ -\beta\mathbb{E}_{(x, y_w, y_l) \sim \mathcal{D}} \bigg[\underbrace{\sigma(\hat{r}_\theta(x, y_l) - \hat{r}_\theta (x, y_w))}_\text{increased emphasis when the reward ranking is inaccurate}\bigg[\underbrace{\nabla_\theta\log \pi(y_w \mid x)}_\text{pushes up $y_w$ probability} - \underbrace{\nabla_\theta\log\pi(y_l \mid x)}_\text{pushes down $y_l$ probability}\bigg]\bigg],
\end{multline*}

where $\hat{r}_\theta(x, y) = \beta \log \frac{\pi_\theta(y \mid x)}{\pi_\text{ref}(y \mid x)}$ defines the implicit reward based on the comparison of the current language model $\pi_\theta$ to a reference model $\pi_\text{ref}$ (see further explanation in Section~\ref{sec:theory}). From an intuitive standpoint, this gradient encourages the model to favor more likely (preferred) completions $y_w$ and reduce the probability assigned to less preferred completions $y_l$. Crucially, each training example's influence is modulated by the degree to which the implicit reward function $\hat{r}_\theta$ mistakenly scores the less favored completion above the preferred one, weighted by $\beta$ to account for the KL constraint's strength. Empirical results in our study demonstrate the necessity of this dynamic weighting: omitting it, as in a simplistic variant of the algorithm, may lead to degeneration of the language model (see Appendix Table~\ref{tab:unlikelihood_generations}).

\textbf{DPO overview.} The standard DPO workflow involves the following steps: (1) For each prompt $x$, generate responses $y_1, y_2 \sim \pi_\text{ref}(\cdot \mid x)$, use human annotation to assign preference labels, and thus create a collection of preference data $\mathcal{D} = \{x^{(i)}, y_w^{(i)}, y_l^{(i)}\}_{i=1}^N$; (2) update the language model $\pi_\theta$ by minimizing $\mathcal{L}_\text{DPO}$ given $\pi_\text{ref}$, the preference data $\mathcal{D}$, and a chosen value of $\beta$. 
In applied settings, it is preferable to leverage publicly available preference datasets rather than producing new samples and collecting fresh human feedback. Given that these datasets are typically drawn from $\pi^\text{SFT}$, we use $\pi^\text{SFT}$ to initialize $\pi_\text{ref}$ whenever possible. In cases where $\pi^\text{SFT}$ is not accessible, $\pi_\text{ref}$ is initialized by maximizing the likelihood on preferred responses ${(x, y_w)}$; formally, ${\pi_\text{ref} = \argmax_{\pi}\mathbb{E}_{x, y_w \sim \mathcal{D}}\left[\log \pi(y_w \mid x)\right]}$. This initialization strategy reduces the distribution mismatch between the inaccessible ground-truth reference and the $\pi_\text{ref}$ actually used in DPO training. Additional implementation and hyperparameter choices are described in Appendix~\ref{app:implementation}.

\section{Theoretical Analysis of DPO}
This section presents a deeper interpretation of the DPO approach, offering formal justification, and drawing connections between the advantages of DPO and the limitations inherent in actor-critic-based RLHF techniques (including PPO~\cite{schulman2017proximal}).

\label{sec:theory}

\subsection{Your Language Model Is Secretly a Reward Model}
DPO achieves both reward modeling and policy learning without requiring separate explicit reward modeling or reinforcement learning stages, instead relying on a unified maximum likelihood objective. In particular, the objective in Eq. \ref{eq:main_eq} corresponds to the Bradley-Terry model, adopting a reward formulation $r^*(x, y) = \beta \log\frac{\pi^*_\theta(y \mid x)}{\pi_\text{ref}(y \mid x)}$; optimization of the parameterized policy $\pi_\theta$ thus aligns with reward model fitting, as per Eq. \ref{eq:reward_model}, after a change of variables. In this section, we elaborate on the theoretical grounding for this reparameterization, demonstrate that it maintains expressivity with respect to the class of learnable reward models, and enables perfect recovery of the optimal policy. We start by introducing an equivalence relation for reward functions.

\begin{definition}
We define two reward functions $r(x, y)$ and $r'(x, y)$ as equivalent if and only if ${r(x, y)-r'(x, y) = f(x)}$ for some function $f$.     
\end{definition}

This relation is straightforwardly seen to satisfy the properties of an equivalence relation, resulting in a partition of reward functions into distinct equivalence classes. We now present two key lemmas:

\begin{lemma}\label{lemma:same_prefrence}
Within the Plackett-Luce framework, which subsumes Bradley-Terry, any pair of reward functions belonging to the same equivalence class gives rise to identical preference distributions.
\end{lemma}

\begin{lemma}\label{lemma:same_policy}
Any two reward functions from a common equivalence class induce the same optimal policy in the context of the constrained RL setting.
\end{lemma}

Since their verification is elementary, we postpone detailed proofs to Appendix \ref{app:lemma1}. The first lemma highlights a classic under-specification phenomenon inherent to the Plackett-Luce model family \cite{plackett1975analysis}. To overcome this ambiguity, it is customary to impose further identifiability restrictions in order to establish properties of the MLE estimates derived from Eq. \ref{eq:reward_model} \cite{bong2022generalized}. The second lemma implies that, from the perspective of the induced optimal policy, all reward functions in a given class are indistinguishable, so for our main objective, it suffices to recover any representative from the optimal equivalence class. The theorem below, demonstrated in Appendix~\ref{app:thm1}, formalizes this statement:

\begin{theorem}\label{thm:main}
    Provided certain mild conditions, every equivalence class of reward functions consistent with Plackett-Luce models (including Bradley-Terry) can be parameterized as ${r(x, y) = \beta \log \frac{\pi(y\mid x)}{\pi_\text{ref}(y\mid x)}}$ for an appropriate model $\pi(y\mid x)$ and reference model $\pi_\text{ref}(y \mid x)$.
\end{theorem}

\begin{sproof}
    For an arbitrary reward function $r(x, y)$, consider the associated optimal policy $\pi_r(y \mid x)$, as described in Eq. \ref{eq:op_policy}. We now show that one can represent a member of the equivalence class of $r$ using the specified reparameterization. Let us introduce the projection operator $f$ as follows:
\begin{equation}
    f(r; \pi_\text{ref}, \beta)(x, y) = r(x, y) - \beta\log\sum_{y}\pi_\text{ref}(y\mid x)\exp\left(\frac{1}{\beta}r(x, y)\right)
\end{equation}
Here, $f$ adjusts $r(x, y)$ by subtracting a normalization constant, specifically, the logarithm of the partition function corresponding to $\pi_r$. Because this normalization term only depends on the prefix $x$, the resulting function $f(r; \pi_\text{ref}, \beta)(x, y)$ resides in the same equivalence class as $r(x, y)$. Replacing $r$ by the expression on the right side of Eq.~\ref{eq:main_eq}, which is applicable for any $r$, we obtain $f(r; \pi_\text{ref}, \beta)(x, y) = \beta \log \frac{\pi_r(y\mid x)}{\pi_\text{ref}(y\mid x)}$. Thus, the operator $f$ yields a representative of the desired form from the equivalence class of $r$, demonstrating that the suggested reparameterization captures the entire set of relevant reward functions.
\end{sproof}

Theorem~\ref{thm:main} can be alternatively interpreted as precisely identifying, within each equivalence class, the particular reward function that the DPO reparameterization chooses; specifically, it selects the reward function that fulfills:

\begin{equation}\label{eq:lag_p}
     \sum_{y}\underbrace{\pi_\text{ref}(y\mid x)\exp\left(\frac{1}{\beta}r(x, y)\right)}_{=\pi(y\mid x)\text{, using Thm.~\ref{thm:main} reparam.}} = 1,
\end{equation}

which ensures that $\pi(y\mid x)$ constitutes a valid probability distribution (i.e., non-negative values summing to unity). Moreover, by referencing Eq.~\ref{eq:op_policy}, we recognize that Eq.~\ref{eq:lag_p} in fact defines the partition function corresponding to the optimal policy associated with the reward function $r(x, y)$. The central innovation in the DPO approach is the imposition of specific restrictions on the otherwise underdetermined Plackett-Luce (and particularly Bradley-Terry) models of preference, thereby maintaining the range of expressible reward functions while rendering the optimal policy from Eq.~\ref{eq:op_policy} analytically computable for all input prompts $x$.

\subsection{Instability of Actor-Critic Algorithms} Our framework can likewise illuminate why standard actor-critic methods for RLHF—such as PPO—may exhibit instability. Focusing on the RL fine-tuning component described in Section~\ref{section:prelims}, we can establish links to the control-as-inference paradigm \cite{levine2018reinforcement} for the constrained RL setting as described by \ref{eq:RL}. Letting $\pi_{\theta}(y\mid x)$ denote a parameterized policy, we seek to minimize $\mathbb{D}_{\text{KL}}[\pi_{\theta}(y|x) \mid \mid \pi^*(y\mid x)]$, where $\pi^*$ is defined as the optimal policy in Eq.~\ref{eq:optimum_model} for reward $r_{\phi}(y, x)$. Straightforward algebra yields the following optimization criterion:

\begin{equation}\label{eq:AC}
    \max_{\pi_{\theta}}\mathbb{E}_{\pi_{\theta}(y\mid x)}\bigg[\underbrace{r_{\phi}(x, y) -\beta\log\sum_{y}\pi_\text{ref}(y\mid x)\exp\left(\frac{1}{\beta}r_{\phi}(x, y)\right)}_{f(r_{\phi}, \pi_\text{ref}, \beta)} - \underbrace{\beta\log\frac{\pi_{\theta}(y\mid x)}{\pi_\text{ref}(y\mid x)}}_{\text{KL}}\bigg]
\end{equation}

The objective under consideration aligns with that of previous studies 
\citep{ziegler2020finetuning, stiennon2022learning, bai2022training, ouyang2022training}, where optimization was conducted using the DPO-equivalent reward for the reward class $r_{\phi}$. Here, the normalization component in $f(r_{\phi}, \pi_\text{ref}, \beta)$ may be viewed as the soft value function corresponding to the reference policy $\pi_\text{ref}$. Although this normalization does not influence the optimal policy, its omission can result in a policy gradient with elevated variance, leading to unstable training dynamics. To mitigate this, one approach is to estimate the normalization through a learned value function; however, this can present challenges in optimization. Another common method, employed in earlier works, involves normalizing rewards using a baseline generated from a human completion—effectively, a single-sample Monte Carlo estimate of the normalization. In contrast, the DPO reparameterization introduces a reward function formulation that obviates the need for such baselines.

\section{Experiments}
This section provides a systematic empirical assessment of DPO’s effectiveness in directly optimizing policies from preference data. Our investigation begins in a controlled environment for text generation, examining the efficiency with which DPO balances reward maximization against KL-divergence minimization relative to the reference policy, in comparison to standard preference optimization algorithms such as PPO. Subsequently, we assess DPO on larger-scale models and more complex RLHF tasks, such as summarization and conversational response. Remarkably, with minimal hyperparameter adjustment, DPO typically achieves comparable or superior outcomes to established baselines, including RLHF with PPO and strategies that select the highest-reward trajectory among $N$ samples based on a learned reward. Prior to detailing our results, we outline our experimental procedures; further methodological specifics can be found in Appendix~\ref{app:exp_details}.

\textbf{Tasks.} The empirical evaluation considers three open-ended text generation tasks. In each setting, algorithms are trained to learn a policy from a collection of preference data $\mathcal{D}=\bigl\{x^{(i)}, y_w^{(i)}, y_l^{(i)}\bigr\}_{i=1}^N$. For the \textbf{controlled sentiment generation} task, $x$ corresponds to an initial segment of a movie review extracted from the IMDb dataset \cite{maas-EtAl:2011:ACL-HLT2011}, with the goal of generating $y$ that exhibits positive sentiment. To facilitate controlled experimentation, preference pairs are created by employing a pre-trained sentiment classifier, ensuring $p(\text{positive}\mid x,y_w)>p(\text{positive}\mid x,y_l)$. Supervised fine-tuning (SFT) is achieved by further training GPT-2-large to convergence using review data from the IMDb training split (additional experimental details are provided in App~\ref{app:sentiment_details}). In the \textbf{summarization} task, $x$ represents a Reddit forum post, with the objective being to generate a summary $y$ capturing the main ideas. Consistent with existing research, experiments utilize the Reddit TL;DR summarization corpus \citep{volske-etal-2017-tl} augmented by human preference labels collected by \citeauthor{stiennon2022learning}. The SFT model here is specifically fine-tuned on forum post summaries authored by humans\footnote{\url{https://huggingface.co/CarperAI/openai_summarize_tldr_sft}} using the TRLX framework \citep{leandro_von_werra_2023_7790115} for reinforcement learning from human feedback (RLHF). The human-labeled preference dataset, collected by \citeauthor{stiennon2022learning}, was drawn from outputs of an SFT model trained in a similar manner but not identical to the one used here. The third task, \textbf{single-turn dialogue}, sets $x$ as a user query ranging from scientific questions to personal advice requests. The policy is trained to generate an informative and appropriate reply $y$. For this, the Anthropic Helpful and Harmless dialogue dataset \citep{bai2022training}, consisting of 170,000 annotated dialogues between humans and an automated assistant, is used. Each dialogue is concluded by two model-generated responses and a human annotation indicating the preferred completion. In this case, since no dedicated SFT model is provided, SFT is carried out by fine-tuning a generic language model exclusively on the preferred responses from the dataset.

\textbf{Evaluation.} We employ two distinct strategies to assess our models. For the controlled sentiment generation setting, where the true reward function (provided by a sentiment classifier) is available, we measure each algorithm’s performance by plotting its trade-off curve between attained reward and KL-divergence relative to the reference policy. This is feasible due to our access to ground-truth rewards. Conversely, in more practical applications where such reward functions are inaccessible, we assess performance based on the algorithm’s \textit{win rate} compared to a baseline policy. Here, GPT-4 serves as a surrogate for human judgments when evaluating summary quality in summarization and response helpfulness in single-turn dialogue. The baseline for summarization is the reference summaries from the test data; for dialogue, we use the human-preferred responses in the dataset. Although some research indicates that LMs can serve as more reliable automated evaluators than conventional metrics \citep{Chen2023ExploringTU}, we perform an independent human study to support our reliance on GPT-4 for evaluations, described in Sec.~\ref{sec:human-judgments}. Our results reveal that GPT-4’s ratings exhibit a high correlation with those of human annotators, with human-GPT-4 agreement matching or exceeding typical inter-annotator reliability.

\textbf{Methods.} Alongside DPO, we examine several pre-existing methods for aligning language model behavior with human preferences. For the summarization task, we use zero-shot prompting with \textbf{GPT-J} \citep{gpt-j}, while in dialogue, we adopt 2-shot prompting using \textbf{Pythia-2.8B} \citep{biderman2023pythia}. Further, we evaluate the \textbf{SFT} model and the \textbf{Preferred-FT} variant, which is fine-tuned using supervised learning on the chosen completions $y_w$, generated either by the SFT model (for sentiment and summarization) or a general LM (for dialogue). We also include the pseudo-supervised \textbf{Unlikelihood} approach~\citep{welleck2019neural}, which encourages the model to increase the probability of $y_w$ while decreasing it for $y_l$, incorporating an adjustable coefficient $\alpha\in[0,1]$ for the unlikelihood term. The study also encompasses \textbf{PPO} \citep{schulman2017proximal}, which is trained with a reward model derived from preference data, as well as \textbf{PPO-GT}, an oracle setup utilizing the authentic reward function in the controlled sentiment context. For sentiment experiments, we compare two PPO-GT variants: a standard implementation \cite{leandro_von_werra_2023_7790115} and a customized version with reward normalization and enhanced hyperparameter tuning for improved outcomes; these modifications are likewise applied to the standard PPO using learned rewards. Additionally, the \textbf{Best of $N$} baseline is included, where $N$ candidate responses are sampled from the SFT model (or Preferred-FT in dialogue), and the best one is selected according to a reward model trained on preference annotations. This method is often highly effective but suffers from limited scalability, as it demands generating $N$ separate completions per query during evaluation.

\subsection{How well can DPO optimize the RLHF objective?}

The KL-constrained reward maximization objective at the core of most RLHF algorithms is designed to maintain a balance between maximizing the reward and ensuring the policy does not diverge excessively from a reference policy. Consequently, when evaluating and comparing these methods, both the attained reward and the corresponding KL divergence must be considered; it is not always beneficial if a method secures a slightly greater reward at the expense of a substantially increased KL divergence. In Figure~\ref{fig:frontier-tldr-main}, the tradeoff between reward and KL divergence is illustrated for several methods within the sentiment task. For each approach, we carry out several independent training sessions, varying the parameter associated with policy conservativeness in each experiment (for PPO, target KL values of $\{3,6,9,12\}$; for $\beta$, values of $\{0.05,0.1,1,5\}$; for unlikelihood, $\alpha$ values of $\{0.05,0.1,0.5,1\}$; and multiple random seeds for preferred-FT). In total, the hyperparameter sweep comprises 22 experiments. At intervals of every 100 training steps, extending through to convergence, we assess each trained policy on a common test prompt set, calculating both the mean reward as determined by the true reward function and the mean sequence-level KL divergence\footnote{Specifically, this refers to summing the KL divergence at each timestep within the sequence.} with respect to the reference policy $\text{KL}\left(\pi\mid \mid \pi_\text{ref}\right)$. The findings show that DPO constructs a remarkably efficient reward-KL tradeoff, delivering the highest rewards with correspondingly low KL divergence. This outcome is especially striking for several reasons. Although both DPO and PPO target the same underlying objective, DPO outperforms PPO in efficiency; the tradeoff achieved by DPO always exceeds that of PPO. Moreover, even when PPO is allowed to utilize the actual ground-truth rewards (PPO-GT), DPO is able to establish a superior reward-KL frontier.

\subsection{Can DPO scale to real preference datasets?}
\label{sec:dpo-real-datasets}
We proceed by assessing the effectiveness of DPO when fine-tuned on tasks involving summarization and single-turn dialogue. It is well-established that standard automatic metrics for summarization, like ROUGE, often show weak alignment with human judgments~\citep{stiennon2022learning}. Previous research demonstrates that leveraging PPO to fine-tune language models based on human preference data can yield more preferable summaries. For our evaluation, we generate samples using various approaches on the test split of the TL;DR summarization dataset, measuring each method’s average win rate relative to reference summaries. Sampling is conducted at temperature settings from 0.0 to 1.0, with corresponding win rates depicted in Figure~\ref{fig:frontier-tldr-main} (right). All methods—including DPO, PPO, and Preferred-FT—begin with the same GPT-J SFT model\footnote{\url{https://huggingface.co/CarperAI/openai_summarize_tldr_sft}}. Notably, DPO achieves a win rate near 61\% at temperature 0.0, outperforming PPO, which reaches approximately 57\% at its optimal sampling temperature. DPO also secures a superior peak win rate compared to the best-of-$N$ baseline. We did not conduct substantial hyperparameter tuning for DPO, specifically for $\beta$, implying that these figures might underrepresent DPO’s best possible performance. Additionally, DPO demonstrates greater stability across different sampling temperatures, unlike PPO, whose outcomes may drop to the level of the original GPT-J model when sampled at higher temperatures. Preferred-FT, on the other hand, offers little to no improvement over the SFT baseline. Section~\ref{sec:human-judgments} presents direct human preference comparisons between DPO and PPO, showing that DPO outputs at a temperature of 0.25 are favored 58\% of the time when contrasted with PPO outputs at the same temperature.

For single-turn dialogue assessment, we examine each method using the subset of the Anthropic HH dataset's test split \citep{bai2022training} that features one exchange between a human and an assistant. In GPT-4-based evaluation, the analysis measures win rates by comparing outputs against the preferred test completions as ground truth. Due to the absence of an established SFT baseline, our experiments begin with the pre-trained Pythia-2.8B model; we employ Preferred-FT to produce a reference model tuned on selected completions to maintain output distribution alignment, followed by DPO training. We also benchmark our approach against the highest scoring output among 128 Preferred-FT completions—since performance plateaus beyond this number (see Appendix Figure~\ref{fig:best-of-n})—and compare with a 2-shot prompted version of the Pythia-2.8B base. Our results indicate that, across optimal temperatures, DPO achieves equivalent or superior outcomes relative to other methods. Additionally, we assess an RLHF model utilizing PPO and trained on the Anthropic HH dataset\footnote{\url{https://huggingface.co/reciprocate/ppo_hh_pythia-6B}} from a prominent repository\footnote{\url{https://github.com/CarperAI/trlx/tree/main/examples/hh}}; however, no prompt or temperature setting yields improvements over the base Pythia-2.8B. Based on these observations, and given that both Best of 128 and PPO optimize identical reward functions (as confirmed by TL;DR experiments), we use Best of 128 as a practical stand-in for PPO performance. Notably, DPO stands out as the only method that is both computationally efficient and capable of surpassing the preferred completions within the Anthropic HH set, rivaling or outperforming the much more resource-intensive Best of 128 approach. Lastly, as depicted in Figure~\ref{fig:dialogue-main}, DPO attains optimal results in comparatively few training steps.

\subsection{Generalization to a new input distribution}

To assess the robustness of PPO and DPO under changes in input distribution, we tested the policies developed from the Reddit TL;DR summarization experiments on a different domain: news articles from the CNN/DailyMail dataset's test split \citep{nallapati-etal-2016-abstractive}. The evaluations used the optimal sampling temperatures previously determined for TL;DR (0 and 0.25), and Table~\ref{tab:ood} reports the findings. The win rate of GPT-4 versus the reference summaries was calculated employing the same GPT-4 (C) prompt as used in the Reddit TL;DR study, with “forum post” replaced by “news article.” In this shifted distribution, DPO substantially outperforms PPO. These results suggest that DPO achieves a level of generalization comparable to PPO, despite DPO not leveraging the extra unlabeled Reddit TL;DR prompts available to PPO.

\subsection{Validating GPT-4 judgments with human judgments}
\label{sec:human-judgments}
To examine the trustworthiness of GPT-4’s evaluation, we conducted a human subject study, utilizing results from the TL;DR summarization experiments and two distinct prompting strategies for GPT-4. The \textbf{GPT-4 (S)} (simple) prompt requests identification of the summary that better captures the essential content of a post. The \textbf{GPT-4 (C)} (concise) prompt, in addition, asks which summary is more concise; this prompt was tested due to our observation that GPT-4—under the simple prompt—shows a preference for lengthier, repetitive summaries compared to human raters. Complete prompt texts are provided in Appendix~\ref{app:prompts}. Three sets of pairwise comparisons were performed: the strongest (DPO, temp. 0.25), the weakest (PPO, temp. 1.0), and an intermediate (SFT, temp. 0.25) system, all measured relative to the best-performing PPO variant (greedy sampling). This selection aimed to sample a range of quality among outputs; for DPO and PPO-1, multiple human ratings were obtained, reflecting resource limitations. Results show that for both prompts, GPT-4’s choices generally correspond to human consensus at similar rates to the agreement found among human annotators themselves. This alignment indicates GPT-4 serves as an adequate stand-in for human evaluators. Notably, the \textbf{GPT-4 (C)} prompt yields win rates that more closely track with human preferences, leading us to adopt it for the primary findings presented in Section~\ref{sec:dpo-real-datasets}. Further information about the human evaluation process, including screenshots of the rating interface and details about volunteer participants, is available in Appendix~\ref{app:human-study}.

\section{Discussion}
Preference-based learning presents an effective and scalable approach for developing capable and aligned language models. In this work, we introduced DPO, a straightforward method for training language models from preferences that eliminates the need for reinforcement learning. Instead of reframing preference learning as a conventional RL problem to leverage existing RL algorithms, DPO constructs a correspondence between language model policies and reward functions, thereby enabling direct optimization for human preferences through a simple cross-entropy loss, all while forgoing reinforcement learning and without sacrificing generality. Notably, DPO achieves results on par with or superior to state-of-the-art RLHF methods, such as those utilizing PPO, with almost no hyperparameter tuning required. This substantially lowers the barrier for leveraging human preferences in language model training.

\textbf{Limitations \& Future Work.} Our findings bring to light several directions for future inquiry. For instance, how robustly does the DPO-trained policy generalize beyond the training distribution compared to models trained on explicit reward signals? Preliminary evidence indicates that DPO models exhibit comparable generalization to those optimized via PPO, but further, more thorough evaluation is necessary. Additional questions include whether using DPO-based self-labeling can effectively utilize prompts without preference labels. Another area to explore is how the phenomenon of reward over-optimization arises in the direct preference optimization paradigm and whether the observed modest performance dip in Figure~\ref{fig:dialogue-main}-right exemplifies this issue. Moreover, although our experiments include models up to 6B parameters, scaling DPO to models with considerably greater capacity remains an intriguing research avenue. In terms of evaluation, we observed that GPT-4-based win rate measurements are influenced by prompt selection; identifying optimal strategies for eliciting accurate assessments from automated evaluators remains an open problem. Lastly, DPO has potential utility across domains beyond aligning language models with human preferences, including the training of generative models for different data modalities.